\subsection*{S2. Strings}

A path holding JSON strings.

\begin{lstlisting}[style=json]
{ "name": "tatami" }
\end{lstlisting}

\options
\begin{enumerate}[label=\textbf{(\alph*)}]
  \item \textbf{OCaml \texttt{string}}, holding the decoded UTF-8 bytes.
        \selected
  \item \textbf{A validated Unicode type}, rejecting or normalising
        ill-formed sequences at ingest.
  \item \textbf{Refine string-encoded scalars} --- recognise dates, UUIDs or
        numbers written as strings and give them narrower types.
\end{enumerate}

\decision{OCaml \texttt{string}. No refinement, no validation beyond what the
parser performs.}

\begin{lstlisting}[style=json]
{ "name": "tatami" }

  -> name : string
\end{lstlisting}

\rationale{Nothing in Phase~1 depends on a finer classification of strings.
The one refinement the project does want --- collapsing a string-encoded
enumeration into an OCaml variant --- is assigned to Phase~2 by the design
note (\S3.3), so introducing it here would duplicate that pass.}

\limitation{OCaml's \texttt{string} is a byte sequence, so nothing guarantees
the contents are well-formed UTF-8 beyond what yojson's decoder enforces on
escape sequences. Option (c) remains available later and is independent of
every other decision here.}
